\documentclass[10pt]{article}
\usepackage[margin=20mm]{geometry}
\usepackage{amsmath,amssymb,mathtools,fontspec,unicode-math,xurl,hyperref}
\IfFontExistsTF{DejaVu Serif}{\setmainfont{DejaVu Serif}}{\setmainfont{Latin Modern Roman}}
\setmathfont{latinmodern-math.otf}
\setlength{\parindent}{0pt}\setlength{\parskip}{6pt}\sloppy
\title{Kernel constants, quantifiers and the proof of Hölder regularity of stationary measures}
\author{AI-assisted technical review}
\date{}
\begin{document}
\maketitle
Abstract: The acquired publisher version of Gorodetski, Kleptsyn and Monakov is examined through its cutoff-energy, transport and contraction arguments. Two local discrepancies are established: a missing dimension factor in an explicit energy constant, and an overbroad monotonicity lemma, for which an exact counterexample is supplied. Their repairs preserve the parameter range and existential conclusions of the main theorems. The normalized Hilbert-space argument, compact-family limit, iteration and moment-boundary example receive positive bounded checks. The report assessment is Minor; scientific adjudication and independent acceptance remain pending.\par
\section*{Identity, assumptions and conclusion scope}
The inspected source is A. Gorodetski, V. Kleptsyn and G. Monakov, Hölder regularity of stationary measures, Inventiones mathematicae 243 (2025), 1037–1077, DOI 10.1007/s00222-025-01389-y (https://doi.org/10.1007/s00222-025-01389-y). It is the acquired publisher version of record, 41 physical pages. Below, PDF pages are one-based and printed page numbers equal PDF page numbers plus 1036. The separately acquired arXiv:2209.12342v1 source is not assumed equivalent to the publisher text.\par
For a smooth closed \(k\)-dimensional Riemannian manifold \(M\), put \(\mathcal L(f)=\max(\operatorname{Lip}(f),\operatorname{Lip}(f^{-1}))\). Theorems 1.1 and 2.3 (PDF pp. 2, 5–6) require a positive moment \(\int\mathcal L(f)^\gamma\,d\mu(f)<\infty\), and absence of a common invariant probability measure. The finite-measure wording of Theorem 2.3 is understood as nonzero finite measures after normalization, consistently with Theorem 1.1. Their conclusion is an upper mass bound \[\nu(B_r(x))\le C r^\beta\] for all stationary probability measures, with uniform positive \(\beta,C\). It does not assert uniqueness, absolute continuity, a lower Ahlfors bound, or exact dimensionality.\par
Theorem 2.6 (PDF p. 6) concerns arbitrary initial probability measures and averaged \(n\)-step images, only at radii \(r>\kappa^n\). Finite-time atomic images are compatible with it. Theorem 2.8 (PDF p. 7) permits arbitrary independent step distributions in a weakly compact family, with a uniform positive moment and the stronger absence of deterministic images: no pair of probability measures \(m,m'\) satisfies \(f_*m=m'\) for almost every \(f\) under any law in that family. These quantifiers and the distinction between stationary and nonstationary hypotheses are essential.\par
This review checks the principal internal dependencies in Sections 2, 4–6 and Appendix A. It does not certify all external results cited in the introductory discussion. Source acquisition and index parsing are recorded as provenance, not as scientific proof verification.\par
\section*{Problem 1: missing dimension factor in the energy ceiling}
\subsection*{Decisive calculation and admissible instance}
Equation reference: eq:correctmax. The error occurs in Lemma 4.7 and the subsequent proof of Corollary 4.8, PDF p. 16 / printed p. 1052. Definition 4.1, PDF p. 13, sets \[\varphi_{\alpha,\varepsilon}(r)=
 \begin{cases}
 r^{-(k+\alpha)/2},&r\ge\varepsilon,\\
 \varepsilon^{-\alpha}r^{-(k-\alpha)/2},&0<r<\varepsilon.
 \end{cases}\] The Euclidean convolution kernel satisfies \(U_{\alpha,\varepsilon}(0)=\int_{\mathbb R^k}
\varphi_{\alpha,\varepsilon}(|x|)^2\,dx\). Writing \(V_k=\operatorname{vol}_{\mathbb R^k}(B_1)\), direct radial integration yields \[\begin{aligned}
 U_{\alpha,\varepsilon}(0)
 &=kV_k\left[
 \varepsilon^{-2\alpha}\int_0^\varepsilon r^{\alpha-1}\,dr
 +\int_\varepsilon^\infty r^{-\alpha-1}\,dr\right]\notag\\
 &=\frac{2kV_k}{\alpha}\varepsilon^{-\alpha}. \end{aligned}\] The source instead defines \(L=2V_k\) and states \(U_{\alpha,\varepsilon}(0)=(L/\alpha)\varepsilon^{-\alpha}\). Its own preceding radial-integration line retains \(k\), so the discrepancy is visible in the actual PDF and is not a text-extraction artefact.\par
Take \(k=2\), \(\alpha=1/4\), \(\varepsilon=1\), within the operative range \(0<\alpha<1/2\). Formula [eq:correctmax] (\#eq:correctmax) gives \(16\pi\), whereas the stated constant gives \(8\pi\). For the admissible probability measure \(\nu=\delta_x\) on any closed two-dimensional Riemannian manifold, \[\mathcal E_{\alpha,\varepsilon}(\delta_x)=U_{\alpha,\varepsilon}(0).\] Thus the specific choice \(C''_\alpha=L/\alpha\) in Corollary 4.8’s proof is too small. The existential assertion that some finite \(C''_\alpha\) works remains true.\par
\subsection*{Repair and actual downstream effect}
Equation reference: eq:holder. The exact repair is \[L=2kV_k,\qquad C''_\alpha=\frac{2kV_k}{\alpha}.\] The proof of Theorem 2.6 explicitly reuses \(C''_\alpha=L/\alpha\) at PDF pp. 24–25 / printed pp. 1060–1061, so this is an affected bound, not merely an unused coefficient. With the corrected value, the same recurrence gives \[\mathcal E_{\alpha,\varepsilon}(\nu_n)
 \le\max\{\lambda^n C''_\alpha\varepsilon^{-\alpha},\widetilde C\}.\] Setting \(\varepsilon=\lambda^{n/\alpha}\) still yields a uniform ceiling and, by Lemma 4.9, \[
 \nu_n(B_r(x))
 \le
 \left(\frac{2^\alpha\max\{2kV_k/\alpha,\widetilde C\}}{c'_\alpha}\right)^{1/2}
 r^{\alpha/2},\qquad r>\lambda^{n/\alpha}.\] The cutoff rate and Hölder exponent are unchanged. Only the explicit finite ceiling and resulting prefactor require enlargement. The initial assessment is Minor, based on this verified dependency boundary. The Dirac example refutes the printed energy bound; it is not asserted to refute a stationary or iterated Hölder bound under the main theorem’s hypotheses.\par
\section*{Problem 2: an exact counterexample to the all-positive-\(\alpha\) lemma}
\subsection*{Source claim and calculation}
Lemma 6.6, PDF p. 27 / printed p. 1063, claims that \(U_{\alpha,\varepsilon}(r)\) is non-increasing in \(r\) for every \(\alpha,\varepsilon>0\). Its proof invokes radial monotonicity of \(\varphi_{\alpha,\varepsilon}\), which need not hold when \(\alpha>k\).\par
Equation reference: eq:uone, eq:utwo. Take \(k=1\), \(\varepsilon=1\), \(\alpha=3\), all allowed by that literal lemma statement. Then \[\varphi(r)=
 \begin{cases}r,&0\le r<1,\\r^{-2},&r\ge1,\end{cases}
 \qquad
 U(r)=\int_{\mathbb R}\varphi(|y|)\varphi(|y-r|)\,dy.\] The functions are nonnegative and square integrable, so the integrals are legitimate and finite. Splitting at the breakpoints and using symmetry gives \[\begin{aligned}
 U(1)
 &=2\int_1^\infty\frac{dt}{t^2(t+1)^2}
   +2\int_0^1\frac{t\,dt}{(t+1)^2}
   +\int_0^1t(1-t)\,dt\notag\\
 &=\frac{13}{6}-2\log2,\\
 U(2)
 &=2\int_1^\infty\frac{dt}{t^2(t+2)^2}
   +2\int_0^1\frac{t\,dt}{(t+2)^2}
   +2\int_0^1\frac{t\,dt}{(2-t)^2}\notag\\
 &=2+\frac32\log3-4\log2.\end{aligned}\] Consequently \[U(2)-U(1)=\frac12\log(27/16)-\frac16
 >\frac{23}{258}>0.\] The strict bound follows from \(\log x>2(x-1)/(x+1)\) for \(x>1\); differentiating the difference gives \((x-1)^2/[x(x+1)^2]>0\). No floating-point assumption is involved. Independent executed arithmetic also returned approximately \(0.7803723\) and \(0.8753297\) for [eq:uone] (\#eq:uone) and [eq:utwo] (\#eq:utwo).\par
\subsection*{Sufficient correction and scope of impact}
Restricting the lemma to \(0<\alpha\le k\) makes both pieces of \(\varphi_{\alpha,\varepsilon}\) non-increasing. The source’s layer-cake argument using intersections of balls then proves the desired radial monotonicity. This is a sufficient repaired range; no maximal range is claimed.\par
Proposition 4.6 and all required main-theorem applications already impose \(0<\alpha<1/2\), contained in that repaired range for positive integer dimension \(k\). Thus the broad auxiliary assertion is false, while the operative small-\(\alpha\) dependency is valid. The initial assessment is Minor. This counterexample is not presented as a counterexample to the main theorem.\par
Lemma 6.8 has a separate broad-\(\alpha\) wording issue in its proof: its displayed uncut-kernel majorant requires \(\alpha<k\). The comparison was verified here in the operative range, but no counterexample to the broader Lemma 6.8 was established. That unverified extension is not conflated with Problem 2.\par
\section*{Verification of the central proof mechanism}
\subsection*{Finite cutoff, energy comparison and positive moments}
For fixed \(\alpha,\varepsilon>0\), the radial exponents of \(\varphi_{\alpha,\varepsilon}^2\), after the volume element, are \(\alpha-1\) near zero and \(-\alpha-1\) at infinity. Hence its Euclidean \(L^2\) norm is finite. The cutoff also gives finite manifold energies for every probability measure, including atoms. This addresses the infinite-energy obstruction expressly identified in Section 3; an inequality of the form \(\mathcal E(\nu)\le\lambda\mathcal E(\nu)+C\) alone would not rule out infinite energy.\par
Tonelli justifies the three-variable kernel representation. The scaling identities in Lemma 6.1 give \[U_{\alpha,\lambda\varepsilon}(\lambda r)
 =\lambda^{-\alpha}U_{\alpha,\varepsilon}(r).\] In the operative range, the convolution limit as \(\varepsilon\downarrow0\) is locally integrable because \(\alpha<k\), and its tail is integrable because \(\alpha>0\). The power-law lower bound and radial monotonicity then yield \[\mathcal E_{\alpha,\varepsilon}(\nu)
 \ge c'_\alpha(2r)^{-\alpha}\nu(B_r(x))^2,\qquad r>\varepsilon,\] with exactly the square-root exponent in [eq:holder] (\#eq:holder).\par
Proposition 4.10 compares the Euclidean interaction energy \(\mathcal E\) with the manifold \(L^2\) energy \(\widetilde{\mathcal E}\) at high energy, uniformly in the cutoff. The normal-coordinate proof is valid after making its near/far split explicit: separated pairs have a uniformly bounded kernel in the small-\(\alpha\) range; close pairs lie in a common normal chart, where the two kernel factors and volume Jacobian are uniformly controlled. The chart tails are bounded using the uncut majorant, without imposing an extra upper restriction on \(\varepsilon\). Finite additive constants are absorbed by the high-energy threshold.\par
The slow-change estimate gives \[\mathcal L(f)^{-\alpha}\mathcal E(\nu)
 \le\mathcal E(f_*\nu)\le\mathcal L(f)^\alpha\mathcal E(\nu).\] If \(\mathbb E\mathcal L^\gamma<C_0\) and \(0<\alpha<\gamma\), then \[C_0^{-\alpha/\gamma}
 <\mathbb E\mathcal L^{-\alpha}
 \le\mathbb E\mathcal L^\alpha
 <C_0^{\alpha/\gamma}.\] The lower bound uses Jensen applied to the reciprocal. These estimates, followed by the energy comparison with additive thresholds, verify the use of Proposition 4.16. They are uniform for the common moment bounds in Theorem 2.8 and do not require continuity of an unbounded moment functional under weak convergence.\par
\subsection*{Normalized variance and transport}
Let \(g_f=\rho_{\alpha,\varepsilon}[f_*\nu]\) and \(g=\mathbb E_\mu g_f=\rho_{\alpha,\varepsilon}[\mu*\nu]\). The cutoff makes the required second moments finite and the normalizing denominators positive. The Hilbert identity gives \[\mathbb E\|g_f-g\|_2^2
 =\mathbb E\widetilde{\mathcal E}(f_*\nu)-\widetilde{\mathcal E}(\mu*\nu).\] Under the noncontraction assumption in Proposition 5.1, this is at most \(2\delta\widetilde{\mathcal E}(\nu)\), while \(\|g\|_2^2\ge(1-\delta)\widetilde{\mathcal E}(\nu)\).\par
For normalized \(h_f=g_f/\|g_f\|_2\) and \(\bar h=g/\|g\|_2\), \[\|h_f-\bar h\|_2\le
 \frac{2\|g_f-g\|_2}{\|g\|_2},\qquad
 \operatorname{TV}(h_f^2\,dy,\bar h^2\,dy)\le\|h_f-\bar h\|_2.\] The source uses half the \(L^1\) difference as total variation. Therefore \[\mathbb E\operatorname{TV}^2\le\frac{8\delta}{1-\delta}\le16\delta
 \quad(\delta\le1/2),\] and Markov gives probability at most \(4\sqrt\delta\) above the threshold \(2\delta^{1/4}\). The Wasserstein estimate follows from \(W_1\le\operatorname{diam}(M)\operatorname{TV}\). This reproduces Lemma 5.3 (PDF pp. 22, 36–37). The repeated \(f_*\nu\) in both printed TV arguments in the later proof is an evident notation slip: the explicitly defined \(\bar h\) and the lemma statement identify the second density as that of \(\mu*\nu\).\par
Proposition 4.21’s approximate equivariance is obtained by comparing the normalized energy measure with the first-coordinate marginal of its three-variable measure. Most mass lies near the diagonal at high energy. The first-coordinate marginal is supported on the original measure’s support, and kernel comparison controls its push-forward for \(\mathcal L(f)\le R\). The triangle inequality and the Lipschitz push-forward bound then give the stated \(W_1\) conclusion. Unspecified constants can absorb the displayed diameter factor, and the small tolerance can be chosen at most one.\par
\subsection*{Compactness, uniformity and iteration}
The order of choices in the contradiction argument is the tolerance, the distortion cutoff \(R\), sufficiently small \(\alpha\), and then a sufficiently large energy threshold. The counterexample measure and cutoff are chosen last. The moment bound supplies the required uniform tail control; for example, if \(\alpha\le\gamma/2\), \[\mathbb E\big[\mathcal L^\alpha 1_{\{\mathcal L>R\}}\big]
 \le C_0R^{-\gamma/2}.\] The resulting normalized measures \(m_j,m'_j\) have subsequential weak limits. In the compact-family case, the laws \(\mu_j\) also converge within the family. The bounded function \[(f,m,m')\longmapsto W_1(f_*m,m')\] is jointly continuous in the \(C^1\) and weak measure topologies on compact \(M\). Its expectation tends to zero by the quantitative exceptional-set estimate. Passing to the limit yields \(\int W_1(f_*m,m')\,d\mu(f)=0\), hence a deterministic image for almost every map, contradicting exactly the assumed condition. No common null set across varying laws is needed.\par
Low-energy control and high-energy contraction combine to give a max recurrence uniform in \(\varepsilon\). Its iteration is valid because \(0<\lambda<1\). Formula [eq:holder] (\#eq:holder) follows with the corrected constant; the energy exponent is renamed after taking its square root. Stationarity then removes the shrinking cutoff for every fixed positive radius.\par
For the weaker stationary hypothesis, Proposition 2.12 (PDF p. 9) supplies a finite block length with no deterministic images. If arbitrarily long deterministic chains existed, their time averages would satisfy \[f_*\bar\nu_\ell-\bar\nu_\ell
 =\frac{\nu_{\ell,\ell}-\nu_{\ell,0}}{\ell}\] on a common full-measure set of maps. Compactness would give a common invariant probability measure. Fixed-length blocking preserves a positive moment by submultiplicativity of \(\mathcal L\). Remainders in \(n\) are absorbed in the arbitrary initial distribution and a finite adjustment of constants and cutoff rate. For a compact family, nested closed sets of laws admitting deterministic block images yield a uniform block length. This also supports the uniformity in Remark 2.4.\par
\section*{Boundary settings and the necessity example}
\subsection*{Manifolds with boundary}
Remarks 2.9 and 5.7 (PDF pp. 7, 25) extend the argument to embeddings of a compact manifold with boundary into itself. An arbitrary such map need not be silently assumed to extend to a diffeomorphism of a chosen ambient closed manifold.\par
A precise adaptation uses the double \(\widetilde M\) and a Lipschitz folding retraction \(P:\widetilde M\to M\). Define the energy densities on the double and use \(\eta(\nu)=P_*\theta(\nu)\) as the probability measure in the final contradiction. The first-coordinate proxy remains supported on \(M\); pushing its comparison with \(\theta\) through \(P\) changes only fixed Lipschitz constants. The kernel ratio uses \(f\) only on \(M\), and the pullback is only on \(f(M)\), so the embedding is sufficient. Total variation contracts under \(P\), preserving the variance estimate. The two compact metrics are uniformly bi-Lipschitz equivalent, so the moment bounds change by a fixed geometric factor. This supports the extension with an explicit adaptation; no global ambient diffeomorphism extension is claimed.\par
\subsection*{Appendix A: checking the moment boundary}
Proposition A.1 (PDF pp. 37–38) uses \[A_n=R_{1/(5n)}
 \begin{pmatrix}2^{n^2}&0\\0&2^{-n^2}\end{pmatrix}
 R_{-1/(5n)},\qquad p_n=2^{-n}.\] These are determinant-one matrices with norm \(2^{n^2}\). Since \(\sum_{n\ge1}n^2 2^{-n}=6\), the logarithmic moment is \(6\log2\). Every positive power moment diverges because \(2^{\gamma n^2-n}\) does not tend to zero. Two matrices with distinct eigendirection pairs, for instance \(A_2,A_3\), exclude a common finite projective invariant set, while each is proximal.\par
With the clockwise projective coordinate, the attracting directions lie in \(I=[0,2/5]\), an invariant cone interval. The attracting directions of \(A_2,A_3\) lie in its interior and their repellers differ. Every point therefore enters that interior under a suitable power of one of these two maps. Compactness supplies a finite cover and a uniform positive finite-block hitting probability; forward invariance and stationarity force all stationary mass into \(I\).\par
On this interval, \[|f'_n(x)|\le q_n:=100\,2^{-2n^2},\qquad
 \nu(B(s_n,q_n^j))\ge p_n^j\] for sufficiently large \(n\). For any integer \(m\ge1\), choose \(n=m+4\). Then \(p_n^m>q_n\), because \(2n^2-mn=m^2+12m+32>7\) and \(2^7>100\). Letting \(j\) grow contradicts any Hölder bound with exponent \(1/m\), and hence any positive exponent. Exact rational computations checked representative values in addition to this analytic proof. The example supports the necessity of the stronger moment assumption; it does not satisfy that assumption and is not a counterexample to the main theorem.\par
\section*{Assessment and limits}
Two local mathematical claims require correction: the explicit dimension-dependent constant and the unrestricted monotonicity quantifier. Both have decisive calculations, source locations and verified repairs. Their core dependencies remain valid after the stated corrections, supporting an overall Minor initial assessment.\par
The broader \(\alpha\)-range of Lemma 6.8 was not resolved, and external introductory literature was not independently re-proved. The completed small-\(\alpha\) checks, boundary adaptation and Appendix A calculations do not imply a universal certification of the paper. No author code is required for these analytic claims, and none was executed. Author adoption of the corrections has not been requested or verified. Scientific decision: pending; human reviews: none; independent acceptance: not performed.\par
\section*{Supporting checks}
New recovery token:\par
No recoverable substantive prior report was present in the assigned snapshot or current input inventory. Their backend identities were not independently verified and their work is not represented as independent acceptance.\par
The primary session actually viewed source pages 1, 2, 5–7, 9, 16, 24–25, 27 and 37–38. Auxiliary source images and exact calculations are separately bound. The accompanying coverage, findings, provenance, compilation and inspection records contain complete source and artifact hashes, commands, logs, dependencies before and after compilation, and all actual report-page image checks.\par
1 A. Gorodetski, V. Kleptsyn and G. Monakov, Hölder regularity of stationary measures, Invent. Math. 243 (2025), 1037–1077. doi:10.1007/s00222-025-01389-y (https://doi.org/10.1007/s00222-025-01389-y).\par
\end{document}
